\documentclass[utf8]{article} %中文文档类
\usepackage[left=2.50cm, right=2.50cm, top=2.50cm, bottom=2.50cm]{geometry} % 页边距
\usepackage{indentfirst} % 首行缩进
\usepackage{fancyhdr,fancybox} % 设置页眉、页脚
\usepackage{ctexcap} % 标题是中文的
\usepackage{helvet} % 用来指定beamer使用的字体
\usepackage{hyperref} % bookmarks 
\usepackage{multicol} % 分栏
\usepackage{cite} % 引用

\usepackage{graphicx} % 图片
\usepackage{subcaption} % 子图
\usepackage{multirow} % 表格
\usepackage{float} % 浮动体
\usepackage{booktabs} % 三线表
\usepackage{longtable} % 长表格
\usepackage{diagbox} % 表头斜线
\usepackage{multirow} % 表格合并
\usepackage{makecell} % 表格换行
\usepackage{enumitem} % 枚举环境宏包
\usepackage{framed} % 边框包

\usepackage{amsmath, amsfonts, amssymb} % 数学公式、符号
\numberwithin{equation}{section}   % 编号从一级标题开始
\usepackage[english]{babel} % 数学公式标准
\usepackage{bm} % 加粗方程字体 
\usepackage{caption} % 标题
\usepackage{color} % 配色
\usepackage[braket]{qcircuit} % 启用狄拉克符号支持, 用于绘制电路图
\usepackage{physics} % 量子力学中的符号加载

\newtheorem{Def}{\hskip 2em {\color{green} 定义}}[subsection]
\newtheorem{Thm}{\hskip 2em {\color{blue} 定理}}[subsection]
\newtheorem{Lem}{\hskip 2em {\color{magenta} 引理}}[subsection]
\newtheorem{Prop}{\hskip 2em {\color{cyan} 命题}}[subsection]
\newtheorem{Note}{\hskip 2em {\color{red} 注}}[subsection]
\newtheorem{Cor}{\hskip 2em {\color{cyan} 推论}}[subsection]
\newtheorem{Eg}{\hskip 2em {\color{black} 例}}[subsection]
\newenvironment{Proof}{{\noindent\it Proof:}\quad}{\hfill $\sharp $\par}
\allowdisplaybreaks

% 页眉页脚设置
\usepackage{fancyhdr}
\usepackage{lastpage} % 获得总页数
\pagestyle{fancy}
\fancyfoot{}
\lhead{\textit{Quantum Computation Basics}}\chead{}\rhead{\textit{\thepage\ of \pageref{LastPage}}}
\lfoot{}\cfoot{}\rfoot{}
\renewcommand{\headrulewidth}{0.4pt}

% 标题设置
\title{\textbf{Model1 量子计算数学基础}}
\date{\today}
\author{\textbf{STU\quad 牛晓铭}}

\hypersetup{
	colorlinks=true,
	linkcolor=black
} % 使目录为黑色字
\begin{document}  
\maketitle
\renewcommand{\contentsname}{} 
\tableofcontents
\thispagestyle{empty}
	
\newpage
\setcounter{page}{1}
\section{量子力学的公理体系(四大假设)}
与经典的 Newton 力学不同, 量子力学并不试图从「直觉」出发, 而是从四条数学公设出发, 推导出一切可观测的结论. 后续的所有算法, 协议和约束都建立在此四大公设基础之上.
\begin{Def}[四大假设概述]
    任何量子力学系统由以下四条公设完全描述: 
    \begin{enumerate}[label=\textbf{公设 \arabic*: }, leftmargin=*]
    \item \textbf{态空间公设} —— 系统状态由 Hilbert 空间中的单位向量描述;
    \item \textbf{演化公设} —— 封闭系统的演化由酉变换(或 Schrödinger 方程)描述;
    \item \textbf{测量公设} —— 测量由一组测量算子描述, 结果满足概率规则;
    \item \textbf{复合系统公设} —— 复合系统的态空间为子系统态空间的张量积.
    \end{enumerate}
\end{Def}

\subsection{公设 1: 态空间公设}

\begin{Thm}
    任意一个孤立的物理系统都与一个称为系统\textbf{状态空间(state space)}的复内积向量空间(即希尔伯特空间)相联系.系统完全由\textbf{状态向量(state vector)}来描述,它是系统空间里的一个单位向量.
\end{Thm}

最重要的量子系统是 \emph{量子比特(quantum bit, qubit)}, 其态空间为 $\mathcal{H} = \mathbb{C}^2$.

\begin{Eg}[qubit]
    一个 qubit 的状态可以写作: 
    \begin{equation*}
        \ket{\psi} = \alpha\ket{0} + \beta\ket{1},
        \quad \alpha, \beta \in \mathbb{C},\quad |\alpha|^2 + |\beta|^2 = 1
    \end{equation*}
    其中 $\{\ket{0}, \ket{1}\}$ 是 $\mathbb{C}^2$ 的一组标准正交基, 
    通常取 $\ket{0} = \begin{bmatrix}1\\0\end{bmatrix}$, $\ket{1} = \begin{bmatrix}0\\1\end{bmatrix}$.
\end{Eg}

与经典比特的区别是本质性的: 经典比特在任何时刻要么是 0 要么是 1, 
而 qubit 可以处于 $\alpha\ket{0} + \beta\ket{1}$ 的\emph{叠加态}.
但注意: 叠加是数学意义上的线性组合, 而非「同时处于 0 和 1」—— 后者是常见的误解.

\subsection{公设 2: 演化公设}

\begin{Thm}
    封闭量子系统的演化可用\textbf{酉变换(unitary transformation)}来描述.也就是说,系统在$t_1$时所处的状态$\ket{\psi}$和在$t_2$时所处的状态$\ket{\psi'}$是通过一个仅于时间$t_1$和$t_2$有关的酉算子$U$联系起来的,即
	\begin{equation}
        \ket{\psi'}=U\ket{\psi}
    \end{equation}
    等价地, 连续时间演化满足 Schrödinger 方程: 
    \begin{equation}
        i\hbar \frac{d}{dt} \ket{\psi(t)} = H(t) \ket{\psi(t)}
    \end{equation}
    在该方程中,$\hbar$是一个物理常数,称为\textbf{普朗克常数(Planck's constant)}.实践中,经常把因子放入,取$\hbar=1$.H是一个称为封闭系统哈密顿量的固定厄米算子.
\end{Thm}

\begin{Note}
    在量子计算中，我们通常考虑离散时间演化的酉门而非连续时间 Schrödinger 方程。
    但 Hamiltonian 模拟需要从 Schrödinger 方程出发构造酉门。
\end{Note}

\subsection{公设 3: 测量公设}

\begin{Thm}
    量子测量由一组测量算子 $\{M_m\}$ 描述, 满足 \emph{完备性关系} $\sum_m M_m^\dagger M_m = I$.
    对状态 $\ket{\psi}$ 进行测量: 
    \begin{enumerate}
        \item 结果 $m$ 出现的概率为 $p(m) = \mel{\psi}{M_m^\dagger M_m}{\psi}$;
        \item 测量后的系统状态坍缩为 $\displaystyle\frac{M_m \ket{\psi}}{\sqrt{p(m)}}$.
    \end{enumerate}
\end{Thm}

完备性方程表达了概率加起来为1的事实:
	\[ 1=\sum_{m}p(m)=\sum_{m}\mel{\psi}{M_m^\dagger M_m}{\psi} \]
这个方程对所有的$\ket{\psi}$都成立,并且等价于完备性方程.但是直接检验完备性方程要容易得多.

\begin{Note}
    测量是量子力学中唯一「非酉」的操作。测量后状态发生了不可逆的改变 ——
    这是量子算法设计中必须考虑的核心约束.
\end{Note}

\begin{Def}
	\textbf{相位因子(phase factor)}是一个复数因子,形式为$e^{i \theta}$.\par 
	考虑状态
	\[ e^{i \theta}\ket{\psi} \]
	其中$\ket{\psi}$是状态向量,$\theta$是一个实数.$e^{i \theta}$称为\textbf{全局相位因子(global phase factor)}.\par 
	考虑状态
	\[ \frac{\ket{0}+\ket{1}}{\sqrt{2}}\quad \frac{\ket{0}-\ket{1}}{\sqrt{2}} \]
	两个状态中$\ket{1}$的振幅分别为$\frac{1}{\sqrt{2}},-\frac{1}{\sqrt{2}}$.若存在实数$\theta$,使得
	\[ a=e^{i \theta}b \]
	则称两个振幅a和b相差一个\textbf{相对相位(relative phase)}.若两个状态的每个振幅都由一个相位因子相联系,则称这两个状态在该基下相差一个相对相位.
\end{Def}
\begin{Note}
	若忽略掉全局相位因子$e^{i \theta}$,状态$e^{i \theta}\ket{\psi}$等于$\ket{\psi}$.对这两种状态的测量统计结果是相同的.假设$M_m$是一个与某个量子测量相关的测量算子,并注意到出现结果m的概率分别为$\mel{\psi}{M_m^\dagger M_m}{\psi}$和$\mel{\psi}{M_m^\dagger M_m}{\psi}=\mel{\psi}{e^{-i \theta}M_m^\dagger M_me^{i \theta}}{\psi}$.因此从观测角度,这两个状态是相同的.
\end{Note}

\subsection{公设 4: 复合系统公设}

\begin{Thm}
    复合物理系统的状态空间是分物理系统的状态空间的张量积.此外,若系统编号为1到n,且系统i的状态被准备为$\ket{\psi_i}$,则整个系统的联合状态是
	\begin{equation}
        \ket{\psi_1}\otimes\ket{\psi_2}\otimes\cdots\otimes\ket{\psi_n}
    \end{equation}
\end{Thm}

张量积结构是量子计算中「指数爆炸」的根源: $n$ 个 qubit 的态空间维度为 $2^n$.
这也是量子计算之所以有潜力的数学基础 —— $n=50$ 时, 态空间维度已经超过 $10^{15}$.

\section{量子比特}
\subsection{单量子比特}
量子比特(qubit)是量子计算的基本信息单元, 其数学形式为 $\mathbb{C}^2$ 中的单位向量: 

\begin{equation}
    \ket{\psi} = \alpha\ket{0} + \beta\ket{1},
    \quad |\alpha|^2 + |\beta|^2 = 1,
    \quad \alpha, \beta \in \mathbb{C}
\end{equation}

\begin{Note}[整体相位]
    态 $\ket{\psi}$ 和 $e^{i\theta}\ket{\psi}$ 在物理上不可区分.
    因此 $\ket{\psi}$ 的自由度实际上是 2(而非 4 —— 因为 $\alpha, \beta$ 各有实部和虚部, 但整体相位消除了一维, 归一化条件又消除了一维).
\end{Note}

\subsection{Bloch 球表示}

\begin{Eg}[单量子比特的 Bloch 球表示]
利用归一化条件 $|\alpha|^2 + |\beta|^2 = 1$, 任意单量子比特态可以写成
\begin{equation}
    \ket{\psi} = e^{i\gamma}\left(\cos\frac{\theta}{2}\,\ket{0} + e^{i\varphi}\sin\frac{\theta}{2}\,\ket{1}\right),
    \qquad \theta \in [0, \pi],\ \varphi, \gamma \in [0, 2\pi).
\end{equation}
忽略不可观测的全局相位 $e^{i\gamma}$ 后, 态由 $(\theta, \varphi)$ 完全确定, 可与单位球面上的点一一对应:
\begin{equation}
    \mathbf{a} = (\sin\theta\cos\varphi,\ \sin\theta\sin\varphi,\ \cos\theta) \in S^2.
\end{equation}
该单位球面称为 \textbf{Bloch 球}. 其中 $\ket{0}$, $\ket{1}$ 分别对应北极和南极;
$\ket{+} = \frac{1}{\sqrt{2}}(\ket{0}+\ket{1})$, $\ket{-} = \frac{1}{\sqrt{2}}(\ket{0}-\ket{1})$ 对应 $x$ 轴上的两个点.
\end{Eg}

\begin{Note}
    Bloch 球表明单量子比特的(物理可区分)态空间是二维曲面 $S^2$.
    量子门的作用可以理解为 Bloch 球上的旋转(见第 \ref{sec:gates} 节).
\end{Note}

\subsection{多量子比特与张量积}

\begin{Thm}
    $n$ 个量子比特组成的复合系统的态空间为
    \begin{equation}
        \mathcal{H} = (\mathbb{C}^2)^{\otimes n} \cong \mathbb{C}^{2^n},
    \end{equation}
    维数为 $2^n$(而非 $2n$). 若各子系统的状态为 $\ket{\psi_i}$, 则复合系统处于\textbf{乘积态(product state)}
    \[ \ket{\psi_1}\otimes\ket{\psi_2}\otimes\cdots\otimes\ket{\psi_n}. \]
\end{Thm}

\begin{Def}[计算基]
    对 $x = x_1\cdots x_n \in \{0,1\}^n$, 记
    \begin{equation}
        \ket{x} := \ket{x_1}\otimes\cdots\otimes\ket{x_n},
    \end{equation}
    集合 $\{\ket{x}: x \in \{0,1\}^n\}$ 称为 $\mathbb{C}^{2^n}$ 的\textbf{计算基(computational basis)}.
    张量积符号常被省略: $\ket{01} \equiv \ket{0,1} \equiv \ket{0}\otimes\ket{1}$, $\ket{0^n} \equiv \ket{0}^{\otimes n}$.
\end{Def}

\begin{Eg}[两量子比特与 Bell 态]
    两量子比特的态空间为 $\mathbb{C}^4$, 标准基为
    \[ \ket{00},\ \ket{01},\ \ket{10},\ \ket{11}. \]
    \textbf{Bell 态}(又称 EPR 对)定义为
    \begin{equation}
        \ket{\Phi^+} = \frac{1}{\sqrt{2}}\big(\ket{00} + \ket{11}\big).
    \end{equation}
    Bell 态无法写成任意乘积态 $\ket{a}\otimes\ket{b}$ 的形式(见第 \ref{sec:constraints} 节),
    这种性质称为\textbf{纠缠(entanglement)}, 是量子计算区别于经典计算的核心资源之一.
\end{Eg}

\begin{Def}[多比特 Pauli 算子]
    作用于第 $i$ 个量子比特（$i = 1, \dots, n$）的 Pauli 算子 $P \in \{X, Y, Z\}$ 定义为
    \begin{equation}
        P_i := I^{\otimes (i-1)} \otimes P \otimes I^{\otimes (n-i)}.
    \end{equation}
    例如两比特系统上 $X_1 = X \otimes I$, $X_2 = I \otimes X$.
\end{Def}

\subsection{密度算子}

\begin{Def}[密度算子]
    若量子系统以概率 $p_i$ 处于纯态 $\ket{\psi_i}$($\sum_i p_i = 1$), 则系统的\textbf{密度算子(density operator)}为
    \begin{equation}
        \rho := \sum_i p_i \ket{\psi_i}\bra{\psi_i}.
    \end{equation}
    纯态 $\ket{\psi}$ 的密度算子为秩 1 矩阵 $\rho = \ket{\psi}\bra{\psi}$.
\end{Def}

\begin{Thm}
    算子 $\rho$ 是某个系综的密度算子, 当且仅当满足
    \begin{enumerate}
        \item $\rho = \rho^\dagger$(厄米性);
        \item $\rho \succeq 0$(半正定性);
        \item $\Tr \rho = 1$(迹归一).
    \end{enumerate}
    若 $\rho^2 = \rho$, 则 $\rho$ 为纯态; 否则称为\textbf{混合态(mixed state)}.
    $n$ 量子比特系统的\textbf{最大混合态}为 $\rho = I/2^n$.
\end{Thm}

\begin{Prop}[密度算子形式下的公设]
    密度算子形式下, 四大公设可以改写为:
    \begin{enumerate}
        \item 演化: $\rho \to U\rho U^\dagger$;
        \item 测量: 结果 $m$ 的概率 $p(m) = \Tr(M_m^\dagger M_m \rho)$, 测量后态为
        $\rho_m = \dfrac{M_m \rho M_m^\dagger}{p(m)}$;
        \item 期望值: $\langle M \rangle = \Tr(M\rho)$.
    \end{enumerate}
    密度算子自动体现全局相位的不可观测性: $\ket{\psi}$ 与 $e^{i\theta}\ket{\psi}$ 对应同一个 $\rho$.
\end{Prop}

\begin{Def}[部分迹与约化密度算子]
    对复合系统 $\mathcal{H}_{AB} = \mathcal{H}_A \otimes \mathcal{H}_B$ 上的算子 $T$,
    定义关于系统 $B$ 的\textbf{部分迹(partial trace)}
    \begin{equation}
        \Tr_B T := \sum_i (I_A \otimes \bra{e_i})\, T\, (I_A \otimes \ket{e_i}),
    \end{equation}
    其中 $\{\ket{e_i}\}$ 是 $\mathcal{H}_B$ 的一组标准正交基.
    系统 $A$ 的\textbf{约化密度算子(reduced density operator)}为 $\rho_A := \Tr_B \rho_{AB}$.
    特别地, 若 $\rho_{AB} = \rho_1 \otimes \rho_2$, 则 $\Tr_B \rho_{AB} = \rho_1$.
    只作用在子系统 $A$ 上的可观测量 $M = M_A \otimes I_B$ 的测量统计只依赖于 $\rho_A$:
    \[ p(m) = \Tr\big(P_m \rho_A\big), \qquad \langle M \rangle = \Tr\big(M_A \rho_A\big). \]
\end{Def}

\begin{Eg}[Bell 态的约化密度算子]
    Bell 态 $\ket{\Phi^+}$ 的密度算子为
    \[ \rho_{AB} = \frac{1}{2}\big(\ket{00}\bra{00} + \ket{00}\bra{11} + \ket{11}\bra{00} + \ket{11}\bra{11}\big). \]
    对系统 $B$ 求部分迹: 对角项 $\Tr_B \ket{00}\bra{00} = \ket{0}\bra{0}$,
    非对角项 $\Tr_B \ket{00}\bra{11} = \braket{0|1}\,\ket{0}\bra{1} = 0$, 于是
    \[ \rho_A = \Tr_B \rho_{AB} = \frac{1}{2}\big(\ket{0}\bra{0} + \ket{1}\bra{1}\big) = \frac{I}{2}. \]
    即复合系统的纯态(Bell 态)在子系统上表现为最大混合态 —— 纠缠的信息完全体现在
    两个子系统之间的关联中, 而不在任何一个子系统内部.
\end{Eg}

\begin{Note}[纯化]
    任意混合态都可以通过引入辅助(ancilla)比特纯化: 若 $\rho = \sum_j p_j \ket{\psi_j}\bra{\psi_j}$,
    则
    \[ \ket{\Psi} = \sum_j \sqrt{p_j}\ket{\psi_j}_A \ket{j}_B \qquad\Longrightarrow\qquad \Tr_B \ket{\Psi}\bra{\Psi} = \rho. \]
\end{Note}

\section{量子逻辑门} \label{sec:gates}

\subsection{单量子比特门}

\begin{Def}[Pauli 门]
    \begin{equation}
        X = \begin{pmatrix} 0 & 1 \\ 1 & 0 \end{pmatrix}, \qquad
        Y = \begin{pmatrix} 0 & -i \\ i & 0 \end{pmatrix}, \qquad
        Z = \begin{pmatrix} 1 & 0 \\ 0 & -1 \end{pmatrix}.
    \end{equation}
    作用效果: $X\ket{0} = \ket{1}$, $X\ket{1} = \ket{0}$(\textbf{比特翻转});
    $Z\ket{0} = \ket{0}$, $Z\ket{1} = -\ket{1}$(\textbf{相位翻转}).
    性质: $X^2 = Y^2 = Z^2 = I$, $XY = iZ$, 且任意两个不同的 Pauli 算子反对易
    ($\{P, Q\} := PQ + QP = 0$). $\{I, X, Y, Z\}$ 构成 $\mathbb{C}^{2\times 2}$ 的一组基.
\end{Def}

\begin{Def}[Hadamard 门、相位门与 T 门]
    \begin{equation}
        H = \frac{1}{\sqrt{2}} \begin{pmatrix} 1 & 1 \\ 1 & -1 \end{pmatrix}, \qquad
        S = \begin{pmatrix} 1 & 0 \\ 0 & i \end{pmatrix}, \qquad
        T = \begin{pmatrix} 1 & 0 \\ 0 & e^{i\pi/4} \end{pmatrix}.
    \end{equation}
    作用: $H\ket{0} = \ket{+}$, $H\ket{1} = \ket{-}$(生成叠加态);
    $S\ket{1} = i\ket{1}$, $T\ket{1} = e^{i\pi/4}\ket{1}$.
    性质: $H^2 = I$, $HXH = Z$, $HZH = X$; $S^2 = Z$, $T^2 = S$.
\end{Def}

\begin{Thm}[旋转门与 Bloch 球旋转]
    对任意 $P \in \{X, Y, Z\}$ 与实数 $\theta$,
    \begin{equation}
        R_P(\theta) := e^{-i\theta P/2} = \cos\frac{\theta}{2}\, I - i\sin\frac{\theta}{2}\, P
    \end{equation}
    是酉算子, 对应 Bloch 球上绕 $P$ 轴旋转角度 $\theta$. 特别地
    \begin{equation}
        R_z(\theta) = \begin{pmatrix} e^{-i\theta/2} & 0 \\ 0 & e^{i\theta/2} \end{pmatrix}.
    \end{equation}
    任何单量子比特酉门都可以写成至多三个旋转门的复合(Euler 分解).
\end{Thm}

\subsection{多量子比特门}

\begin{Def}[CNOT 门]
    CNOT(受控非门)的作用为
    \begin{equation}
        \text{CNOT}\ket{a}\ket{b} = \ket{a}\ket{a \oplus b},
        \qquad a \oplus b = (a + b) \bmod 2 \ \text{为异或(XOR)运算},
    \end{equation}
    其中 $\ket{a}$ 称为\textbf{控制比特}, $\ket{b}$ 称为\textbf{目标比特}.
    在计算基 $\{\ket{00}, \ket{01}, \ket{10}, \ket{11}\}$ 下的矩阵表示为
    \begin{equation}
        \text{CNOT} = \begin{pmatrix} 1 & 0 & 0 & 0 \\ 0 & 1 & 0 & 0 \\ 0 & 0 & 0 & 1 \\ 0 & 0 & 1 & 0 \end{pmatrix}.
    \end{equation}
    电路表示为
    \begin{equation*}
    \Qcircuit @C=1em @R=.8em {
        \lstick{\ket{a}} & \ctrl{1} & \rstick{\ket{a}} \qw \\
        \lstick{\ket{b}} & \targ & \rstick{\ket{a \oplus b}} \qw
    }
    \end{equation*}
\end{Def}

\begin{Def}[受控酉门]
    对任意酉算子 $U$, \textbf{受控-$U$ 门}定义为
    \begin{equation}
        CU = \ket{0}\bra{0} \otimes I + \ket{1}\bra{1} \otimes U.
    \end{equation}
    CNOT 即 $U = X$ 的情形. 空圈表示以 $\ket{0}$ 为控制态, 可用于实现量子版的条件语句(量子 if).
\end{Def}

\begin{Def}[SWAP 门与 Toffoli 门]
    SWAP 门交换两个比特的状态: $\text{SWAP}\ket{a}\ket{b} = \ket{b}\ket{a}$, 电路为
    \begin{equation*}
    \Qcircuit @C=1em @R=.8em {
        \lstick{\ket{a}} & \qswap & \qw & \rstick{\ket{b}} \qw \\
        \lstick{\ket{b}} & \qswap \qwx & \qw & \rstick{\ket{a}} \qw
    }
    \end{equation*}
    (SWAP 门由两条导线上、下对齐的两个 $\times$ 符号及连接它们的竖线构成.)
    \textbf{Toffoli 门}(受控受控非门, CCNOT)的作用为
    \begin{equation}
        \text{Toffoli}\ket{a}\ket{b}\ket{c} = \ket{a}\ket{b}\ket{c \oplus (a \wedge b)},
    \end{equation}
    电路为
    \begin{equation*}
    \Qcircuit @C=1em @R=.8em {
        \lstick{\ket{a}} & \ctrl{1} & \rstick{\ket{a}} \qw \\
        \lstick{\ket{b}} & \ctrl{1} & \rstick{\ket{b}} \qw \\
        \lstick{\ket{c}} & \targ & \rstick{\ket{c \oplus (a \wedge b)}} \qw
    }
    \end{equation*}
\end{Def}

\subsection{通用门集}

\begin{Def}
    一个门集 $\mathcal{S}$ 是\textbf{通用(universal)}的, 若对任意 $n$ 比特酉算子 $U$ 与任意精度
    $\varepsilon > 0$, 存在 $\mathcal{S}$ 中的门序列 $U_1, \dots, U_m$ 使得
    $\norm{U - U_m\cdots U_1} \le \varepsilon$(范数取算子范数 $\norm{A} = \sup_{\norm{v}=1}\norm{Av}$).
    注意通用性允许近似实现 —— 精确实现一般需要连续无穷多门.
\end{Def}

\begin{Thm}[通用门集的例子]
    $\{H, T, \text{CNOT}\}$ 与 $\{H, \text{Toffoli}\}$ 都是通用门集.
\end{Thm}

\begin{Thm}[Solovay--Kitaev 定理]
    设 $\mathcal{S}$, $\mathcal{T}$ 是两个对逆封闭的通用门集. 则用 $\mathcal{T}$ 中的门模拟
    $\mathcal{S}$ 中 $m$ 个门构成的电路到精度 $\varepsilon$, 只需 $O\big(m \cdot \mathrm{polylog}(m/\varepsilon)\big)$ 个门;
    且存在经典算法在 $O\big(m \cdot \mathrm{polylog}(m/\varepsilon)\big)$ 时间内找到该电路.
    换言之, 通用门集的选择对量子算法的复杂度只造成多项式对数级别的差别.
\end{Thm}

\begin{Def}[Pauli 群与 Clifford 群]
    $n$ 比特 \textbf{Pauli 群} $\mathcal{P}_n$ 由
    $i^k\, P_1 \otimes \cdots \otimes P_n$($P_j \in \{I, X, Y, Z\}$, $k \in \{0, 1, 2, 3\}$)
    组成, 共 $4^{n+1}$ 个元素.
    \textbf{Clifford 群} $\mathcal{C}_n$ 是共轭作用下保持 Pauli 群不变的酉算子集合:
    \begin{equation}
        \mathcal{C}_n = \big\{C \in U(2^n) : C P C^\dagger \in \mathcal{P}_n,\ \forall P \in \mathcal{P}_n\big\},
    \end{equation}
    可由 $\{H, S, \text{CNOT}\}$ 生成.
\end{Def}

\begin{Note}[Gottesman--Knill 定理]
    仅含 Clifford 门、初态与测量都在计算基上的量子电路, 可以在经典计算机上高效模拟.
    因此通用量子计算必须引入非 Clifford 门; 容错量子计算中常用的通用门集是 \textbf{Clifford+T}.
\end{Note}

\subsection{量子电路与可逆计算}

\begin{Def}[电路约定]
    量子电路中时间从左向右流动, 每条线(wire)代表一个量子比特, 自上而下依次为第 $0, 1, \dots$ 个量子比特.
    若干量子比特组成\textbf{寄存器(register)}: 承载主要计算状态的为\textbf{系统寄存器},
    辅助计算过程、最终需要还原的为\textbf{工作寄存器(working register)},
    用于控制其它门的为\textbf{控制寄存器}.
\end{Def}

\begin{Prop}[经典计算的量子化]
    量子计算机至少与经典计算机一样强大: 任何经典电路都可以高效地用量子电路模拟.
    关键在于把不可逆的经典门\textbf{可逆化}:
    \begin{equation}
        (x, y) \;\longrightarrow\; (x,\ y \oplus f(x)),
    \end{equation}
    存储所有中间计算步骤(即保存 $x$), 再用酉算子实现该可逆映射.
\end{Prop}

\begin{Def}[Oracle]
    \textbf{Oracle(预言机)}是模拟布尔函数 $f: \{0,1\}^n \to \{0,1\}^m$ 的黑盒酉算子:
    \begin{equation}
        U_f \ket{x}\ket{y} = \ket{x}\ket{y \oplus f(x)}.
    \end{equation}
    若 $f$ 输出单比特, 也可用\textbf{相位 oracle} $U_f \ket{x} = (-1)^{f(x)}\ket{x}$,
    这称为\textbf{相位反冲(phase kickback)}: $f(x)$ 的值以相位因子的形式返回,
    可以同时作用于叠加态中的所有基态.
\end{Def}

\begin{Note}[Uncomputation]
    经典计算中存储的中间结果在量子计算中不能直接丢弃: 残留的垃圾寄存器与系统寄存器之间存在纠缠,
    会干扰后续计算. 处理方法是对可逆电路再施加一遍\textbf{逆运算}(uncomputation),
    把工作寄存器恢复为 $\ket{0^m}$, 使其不再与系统纠缠, 从而可以复用.
    这是量子算法设计中最重要的工程原则之一.
\end{Note}

\section{量子测量} \label{sec:measurement}

\subsection{投影测量与可观测量}

\begin{Thm}[Born 规则]
    量子\textbf{可观测量(observable)}由态空间上的厄米算子 $M$ 描述, 其谱分解为
    \begin{equation}
        M = \sum_m m P_m,
    \end{equation}
    其中 $m \in \mathbb{R}$ 为特征值(因为 $M$ 厄米, 特征值必为实数, 从而可以解释为测量结果),
    $P_m$ 为到特征子空间的投影算子($P_m^2 = P_m$, $P_m P_{m'} = 0$).
    对处于纯态 $\ket{\psi}$ 的系统测量 $M$:
    \begin{enumerate}
        \item 测量结果必定是某个特征值 $m$, 概率为
        \begin{equation}
            p_m = \mel{\psi}{P_m}{\psi};
        \end{equation}
        \item 测量后系统坍缩到对应特征子空间, 即
        \begin{equation}
            \ket{\psi} \;\longrightarrow\; \frac{P_m \ket{\psi}}{\sqrt{p_m}}.
        \end{equation}
    \end{enumerate}
\end{Thm}

投影算子满足\textbf{完备性关系}(单位分解) $\sum_m P_m = I$, 从而
\[ \sum_m p_m = \sum_m \mel{\psi}{P_m}{\psi} = \braket{\psi|\psi} = 1, \]
即 $\{p_m\}$ 构成概率分布.

\begin{Cor}[期望值]
    测量结果的期望值为
    \begin{equation}
        \langle M \rangle := \sum_m m\, p_m = \sum_m m \mel{\psi}{P_m}{\psi}
        = \mel{\psi}{\left(\sum_m m P_m\right)}{\psi} = \mel{\psi}{M}{\psi}.
    \end{equation}
\end{Cor}

\begin{Eg}[测量 Pauli-$X$]
    $X$ 的特征分解为 $X = \ket{+}\bra{+} - \ket{-}\bra{-}$, 特征值为 $\pm 1$,
    其中 $\ket{\pm} = \frac{1}{\sqrt{2}}(\ket{0} \pm \ket{1})$.
    对 $\ket{\psi} = \ket{0} = \frac{1}{\sqrt{2}}(\ket{+} + \ket{-})$ 测量 $X$:
    \[ p_{+} = p_{-} = \frac{1}{2}, \qquad \langle X \rangle = 0. \]
\end{Eg}

\subsection{一般测量算子与 POVM}

\begin{Thm}
    一般地, 量子测量由一组\textbf{测量算子(measurement operators)} $\{M_m\}$ 描述,
    满足\textbf{完备性关系}
    \begin{equation}
        \sum_m M_m^\dagger M_m = I.
    \end{equation}
    对状态 $\ket{\psi}$:
    \begin{enumerate}
        \item 结果 $m$ 出现的概率为 $p(m) = \mel{\psi}{M_m^\dagger M_m}{\psi}$;
        \item 测量后的状态坍缩为 $\dfrac{M_m \ket{\psi}}{\sqrt{p(m)}}$.
    \end{enumerate}
    投影测量是 $M_m = P_m$ 的特殊情形. 完备性关系等价于 $\sum_m p(m) = 1$ 对所有 $\ket{\psi}$ 成立.
\end{Thm}

\begin{Def}[POVM]
    令 $E_m := M_m^\dagger M_m$, 则 $E_m \succeq 0$ 且 $\sum_m E_m = I$.
    算子族 $\{E_m\}$ 称为\textbf{正算子值测度(positive operator-valued measure, POVM)}.
    对密度算子 $\rho$, 测量结果 $m$ 的概率为 $p(m) = \Tr(E_m \rho)$.
    与投影测量不同, POVM 的元素 $E_m$ 不必是投影算子($E_m^2 \neq E_m$), 也不必相互正交.
\end{Def}

\begin{Thm}[Naimark 延拓定理]
    任何 POVM $\{E_m\}$ 都可以通过在更大的 Hilbert 空间上的投影测量实现:
    存在辅助空间 $\mathcal{H}_B$、纯态 $\ket{0}_B$ 与投影测量 $\{P_m\}$, 使得对任意 $\rho$,
    \[ \Tr(E_m \rho) = \Tr\!\big(P_m (\rho \otimes \ket{0}\bra{0}_B)\big). \]
    因此在量子算法设计中, 通常只需考虑投影测量.
\end{Thm}

\subsection{测量的两个基本原理}

\begin{Prop}[测量延迟原理]
    测量可以从量子电路的中间阶段移到电路的末尾: 若中间测量结果被用来对后续门作经典控制,
    则该经典控制总可以等价地改写为量子控制, 并把测量推迟到电路末端.
    推迟时需要额外的辅助比特来存储中间测量信息.
\end{Prop}

\begin{Eg}[测量延迟需要额外比特]
    电路 $\ket{0} \xrightarrow{H} (\text{测量}) \xrightarrow{H}$ 以概率 $1/2$ 输出 $0$ 或 $1$;
    若简单地删去中间测量并把两个 $H$ 连续作用, 则 $\ket{0} \xrightarrow{HH} \ket{0}$ 恒输出 $0$.
    正确的做法是用 CNOT 把测量结果拷贝到辅助比特上, 保留中间信息后再推迟测量.
\end{Eg}

\begin{Prop}[隐式测量原理]
    在电路末尾, 未被测量的量子比特可以假设为已被测量: 由约化密度算子的性质可知,
    测量子系统 $A$ 的统计结果不依赖于子系统 $B$ 是否被测量.
\end{Prop}

\subsection{期望值的估计: Hadamard 检验}

\begin{Eg}[实部]
    Hadamard 检验用于估计酉算子 $U$ 在态 $\ket{\psi}$ 上的期望值 $\mel{\psi}{U}{\psi}$:
    \begin{equation*}
    \Qcircuit @C=1em @R=.8em {
        \lstick{\ket{0}} & \gate{H} & \ctrl{1} & \gate{H} & \meter & \rstick{0/1} \qw \\
        \lstick{\ket{\psi}} & \qw & \gate{U} & \qw & \qw & \qw
    }
    \end{equation*}
    电路将 $\ket{0}\ket{\psi}$ 变换为
    \[ \frac{1}{2}\ket{0}\big(\ket{\psi} + U\ket{\psi}\big) + \frac{1}{2}\ket{1}\big(\ket{\psi} - U\ket{\psi}\big), \]
    因此测量第一个量子比特输出 $0$ 的概率为
    \begin{equation}
        p(0) = \frac{1}{2}\Big(1 + \Re \mel{\psi}{U}{\psi}\Big).
    \end{equation}
\end{Eg}

\begin{Eg}[虚部]
    在控制比特上、受控门之前施加相位门的共轭 $S^\dagger = \mathrm{diag}(1, -i)$:
    \begin{equation*}
    \Qcircuit @C=1em @R=.8em {
        \lstick{\ket{0}} & \gate{H} & \gate{S^\dagger} & \ctrl{1} & \gate{H} & \meter & \rstick{0/1} \qw \\
        \lstick{\ket{\psi}} & \qw & \qw & \gate{U} & \qw & \qw & \qw
    }
    \end{equation*}
    此时电路将 $\ket{0}\ket{\psi}$ 变换为
    \[ \frac{1}{2}\ket{0}\big(\ket{\psi} - iU\ket{\psi}\big) + \frac{1}{2}\ket{1}\big(\ket{\psi} + iU\ket{\psi}\big), \]
    输出 $0$ 的概率为
    \begin{equation}
        p(0) = \frac{1}{2}\Big(1 + \Im \mel{\psi}{U}{\psi}\Big).
    \end{equation}
    结合实部、虚部两个电路, 即可完整恢复复数期望值
    $\mel{\psi}{U}{\psi} = \Re\mel{\psi}{U}{\psi} + i\,\Im\mel{\psi}{U}{\psi}$.
\end{Eg}

\begin{Eg}[SWAP 检验]
    SWAP 检验是 Hadamard 检验在 $U$ 取 SWAP 门时的特例, 用于估计两个量子态
    $\ket{\psi}$, $\ket{\varphi}$ 的重叠度:
    \begin{equation*}
    \Qcircuit @C=1em @R=.8em {
        \lstick{\ket{0}} & \gate{H} & \ctrl{1} & \gate{H} & \meter & \rstick{0/1} \qw \\
        \lstick{\ket{\psi}} & \qw & \qswap & \qw & \qw & \qw \\
        \lstick{\ket{\varphi}} & \qw & \qswap \qwx & \qw & \qw & \qw
    }
    \end{equation*}
    其中受控 SWAP 门(又称 Fredkin 门)由控制比特同时控制两个目标比特上的交换操作.
    测量第一个量子比特输出 $0$ 的概率为
    \begin{equation}
        p(0) = \frac{1}{2}\Big(1 + |\braket{\psi|\varphi}|^2\Big).
    \end{equation}
\end{Eg}

\begin{Note}
    SWAP 检验只能估计重叠度的模平方 $|\braket{\psi|\varphi}|^2$ —— 相位信息在测量中被丢弃.
    若需要复数 $\braket{\psi|\varphi}$ 本身, 可用更精细的相位估计类方法(如 swap test 的推广).
\end{Note}

\subsection{测量与采样误差}

\begin{Prop}
    用测量估计期望值 $\langle M \rangle = \Tr(M\rho)$ 到\textbf{加性(additive)}精度 $\varepsilon$,
    所需样本数满足
    \begin{equation}
        N \gtrsim \frac{\mathrm{Var}(M)}{\varepsilon^2},
        \qquad \mathrm{Var}(M) := \langle M^2 \rangle - \langle M \rangle^2.
    \end{equation}
    例如估计单比特输出 $1$ 的概率 $p$(取 $M = \ket{1}\bra{1}$, 则 $\mathrm{Var} = p(1-p)$),
    需要 $N \gtrsim p(1-p)/\varepsilon^2$ 次采样.
\end{Prop}

\begin{Note}
    若 $p$ 接近 $0$ 或 $1$, 加性精度下采样次数很少; 但若要求\textbf{乘性(multiplicative)}精度
    (相对误差 $\varepsilon$), 则 $N \gtrsim \dfrac{1-p}{p\,\varepsilon^2}$, 当 $p \to 0$ 时代价急剧增大.
    这一观测是振幅放大(amplitude amplification)等技巧的出发点.
\end{Note}

\section{量子力学基本约束} \label{sec:constraints}

\subsection{不可克隆定理}

\begin{Thm}[不可克隆定理]
    不存在酉算子 $U$ 能够对任意未知量子态 $\ket{\psi}$ 完成复制, 即不存在固定的辅助态 $\ket{s}$ 使得
    \begin{equation}
        U \ket{\psi}\ket{s} = \ket{\psi}\ket{\psi}, \qquad \forall\, \ket{\psi}.
    \end{equation}
\end{Thm}

\begin{Proof}
    假设这样的 $U$ 存在. 取两个态 $\ket{\psi_1}$, $\ket{\psi_2}$ 满足 $0 < |\braket{\psi_1|\psi_2}| < 1$, 则
    \[ \braket{\psi_1|\psi_2} = \bra{\psi_1}\bra{s}\, U^\dagger U\, \ket{\psi_2}\ket{s}
       = \braket{\psi_1|\psi_2}^2, \]
    迫使 $\braket{\psi_1|\psi_2} \in \{0, 1\}$, 与假设矛盾.
\end{Proof}

\begin{Note}[两个例外]
    不可克隆定理有两个重要例外:
    \begin{enumerate}
        \item \textbf{已知态的制备}: 若 $\ket{\psi} = U_\psi \ket{s}$ 且 $U_\psi$ 已知, 则
        $(I \otimes U_\psi)\ket{\psi}\ket{s} = \ket{\psi}\ket{\psi}$;
        \item \textbf{计算基上的经典信息}: $\text{CNOT}\ket{x}\ket{0} = \ket{x}\ket{x}$,
        $x \in \{0,1\}$, 即经典比特可以复制.
    \end{enumerate}
    事实上, 酉算子可以复制一组\textbf{相互正交}的态, 而一般叠加态不能被复制.
\end{Note}

\begin{Note}[对科学计算的启示]
    经典迭代算法(如求解线性方程组的迭代法)需要不断存储中间变量 $y \leftarrow x$,
    这在量子计算中一般无法直接实现. 这迫使量子算法采用与经典算法本质不同的结构,
    例如用振幅或相位编码信息、用 uncomputation 释放寄存器.
\end{Note}

\subsection{不可删除定理}

\begin{Thm}[不可删除定理]
    不存在酉算子 $U$ 将任意未知量子态重置为 $\ket{0^n}$:
    \begin{equation}
        U \ket{0^n}\ket{\psi} = \ket{0^n}\ket{0^n}, \qquad \forall\, \ket{\psi}.
    \end{equation}
\end{Thm}

\begin{Proof}
    若存在, 取正交态 $\ket{\psi_1}$, $\ket{\psi_2}$, 则由酉性
    \[ 0 = \braket{\psi_1|\psi_2} = \bra{0^n}\bra{\psi_1}\, U^\dagger U\, \ket{0^n}\ket{\psi_2}
       = \braket{0^n|0^n}^2 = 1, \]
    矛盾.
\end{Proof}

\begin{Note}
    更强的版本: 给定任意未知态的\textbf{两份拷贝}, 也不存在酉算子删除其中一份而保留另一份.
    不可删除定理与不可克隆定理都是量子力学线性性的直接推论.
\end{Note}

\subsection{测量坍缩的不可逆性}

\begin{Note}
    测量是量子力学中唯一\textbf{非酉}、不可逆的操作: 坍缩过程丢失了叠加信息,
    无法通过酉演化撤销. 这对算法设计构成核心约束:
    \begin{enumerate}
        \item 测量结果本质上是概率性的, 量子算法通常只能保证以足够高的概率给出正确答案,
        或通过多次运行估计期望值;
        \item 应尽量把测量推迟到电路末端(测量延迟原理), 中间过程保持酉性;
        \item 中间测量结果的经典控制等价于量子控制, 但需要辅助比特保存信息.
    \end{enumerate}
\end{Note}

\subsection{纠缠: 超越乘积态的约束与资源}

\begin{Eg}[Bell 态不是乘积态]
    设 $\ket{\Phi^+} = \ket{a}\otimes\ket{b}$ 是乘积态, 其中
    $\ket{a} = \alpha_0\ket{0} + \alpha_1\ket{1}$, $\ket{b} = \beta_0\ket{0} + \beta_1\ket{1}$, 则
    \[ \ket{\Phi^+} = \alpha_0\beta_0\ket{00} + \alpha_0\beta_1\ket{01} + \alpha_1\beta_0\ket{10} + \alpha_1\beta_1\ket{11}. \]
    与定义 $\ket{\Phi^+} = \frac{1}{\sqrt{2}}(\ket{00} + \ket{11})$ 比较系数:
    $\alpha_0\beta_1 = \alpha_1\beta_0 = 0$ 且 $\alpha_0\beta_0 = \alpha_1\beta_1 = 1/\sqrt{2} \neq 0$,
    即 $\alpha_0, \alpha_1, \beta_0, \beta_1$ 必须全部非零, 矛盾.
\end{Eg}

\begin{Note}
    复合系统的态空间 $\mathbb{C}^{2^n}$ 远大于乘积态构成的子集 —— 绝大多数量子态都是纠缠态.
    纠缠既是量子计算的资源(量子纠错、量子通信、相位估计等依赖纠缠), 也带来约束:
    例如 Bell 态的约化密度算子是最大混合态, 说明子系统中不含任何可被局部观测到的信息.
\end{Note}

\section{量子算法的评价框架}

\subsection{量子算法的一般结构}

\begin{Def}
    一个典型量子算法的组成部分为:
    \begin{enumerate}
        \item 初态: 通常为计算基态 $\ket{0^n}$(或由已知制备程序 $U_x \ket{0^n}$ 得到的输入态 $\ket{x}$);
        \item 酉电路: 由通用门集中的单、双比特门按时间顺序复合而成;
        \item 末端测量: 在计算基上测量(部分)量子比特, 读取结果.
    \end{enumerate}
    系统中的量子比特分为系统寄存器与辅助(ancilla)寄存器; 辅助寄存器中经 uncomputation
    可复位复用的称为工作寄存器. 由于酉性要求, 中间测量应推迟到末端(见第 \ref{sec:measurement} 节).
\end{Def}

\subsection{复杂度度量}

\begin{Def}[高效算法]
    设 $n$ 为输入规模（编码输入所需的量子比特数）. 若电路中的门数为 $O(\mathrm{poly}(n))$ 量级,
    则称该算法为\textbf{高效的}. 由于测量结果具有概率性, 通常要求算法以足够高的概率
    (如 $p > 2/3$)输出正确答案; 通过重复运行并取多数投票(Chernoff 界), 可以把错误概率压到任意小.
\end{Def}

\begin{Def}[三种复杂度]
    \begin{enumerate}
        \item \textbf{门复杂度(gate complexity)}: 电路中量子门的总数;
        \item \textbf{查询复杂度(query complexity)}: 对黑盒 oracle 的调用次数
        (例如 Grover 搜索对 oracle 的查询下界为 $\Omega(\sqrt{N})$);
        \item \textbf{电路深度(circuit depth)}: 从输入到输出的任意路径上的最大门数,
        近似对应真实运行时间(墙钟时间).
    \end{enumerate}
\end{Def}

\begin{Note}
    查询复杂度隐藏了 oracle 本身的实现代价, 门复杂度与查询复杂度之间可能存在巨大差距:
    有些 oracle(如求逆矩阵对应的 oracle)实现起来极难.
    因此分析算法时, 应同时确认其它组件不会主导总代价.
    另外, 量子态只能保持有限时间(相干时间), 电路深度应尽量小, 即使这意味着需要重复运行更多次.
\end{Note}

\subsection{误差分析: 混合论证与误差线性增长}

\begin{Prop}[混合论证]
    设 $U_i$, $\tilde U_i$ 为酉算子且 $\norm{U_i - \tilde U_i} \le \varepsilon$($i = 1, \dots, K$), 则
    \begin{equation}
        \norm{U_K\cdots U_1 - \tilde U_K \cdots \tilde U_1} \le K \varepsilon.
    \end{equation}
    即若每个局部酉门都精确到 $\varepsilon$, 全局误差至多线性增长为 $K\varepsilon$:
    要实现总精度 $\varepsilon_{\mathrm{tot}}$, 每个门的精度需达到 $\varepsilon_{\mathrm{tot}}/K$ 量级.
\end{Prop}

\begin{Proof}
    使用望远镜级数展开:
    \begin{equation*}
        U_K\cdots U_1 - \tilde U_K\cdots\tilde U_1
        = \sum_{i=1}^{K} U_K\cdots U_{i+1}\, (U_i - \tilde U_i)\, \tilde U_{i-1}\cdots \tilde U_1,
    \end{equation*}
    对每一项取算子范数, 由酉算子的范数不变性($\norm{U A} = \norm{A}$)与三角不等式即得结论.
\end{Proof}

\begin{Prop}[Hamiltonian 模拟的 Duhamel 原理]
    设 $i\partial_t U = HU$, $i\partial_t \tilde U = H\tilde U + B(t)$, $U(0) = \tilde U(0) = I$,
    其中 $H$ 厄米, $B(t)$ 任意, 则
    \begin{equation}
        \norm{U(t) - \tilde U(t)} \le \int_0^t \norm{B(s)}\, ds.
    \end{equation}
    这是混合论证的连续时间版本(常数变易法), 是分析 Trotter 分裂等 Hamiltonian 模拟误差的基础.
\end{Prop}

\subsection{容错计算与阈值定理}

\begin{Thm}[阈值定理]
    若单个量子门的噪声低于某个常数阈值(约为 $10^{-4}$ 量级, 依赖具体的错误模型与纠错码),
    则可以通过量子纠错以任意期望精度、仅付出多项式对数级别的开销, 高效执行任意大的量子计算.
\end{Thm}

\begin{Note}
    本系列讲义始终假设容错协议已经实现, 因此只需关心两类误差:
    数学层面的\textbf{近似误差}(如 Trotter 分裂、函数多项式逼近带来的误差),
    以及读出阶段的\textbf{蒙特卡洛(采样)误差}(由于测量结果本质随机).
    后续章节对算法复杂度的评价都建立在这个框架之上.
\end{Note}


\end{document}